\chapter{HotSpot, BEAM e GraalVM}
\label{cap:comparacao-vms}

\proximasaulas

As Partes~\ref{parte:java-jvm} e~\ref{parte:quarkus} olharam uma VM só: a HotSpot, com e sem Quarkus.
Este capítulo põe outras duas ao lado dela, para que as escolhas da
HotSpot apareçam como escolhas, e não como a única forma de fazer uma VM.

\begin{description}
\item[HotSpot] A JVM do OpenJDK, que a Parte~\ref{parte:java-jvm} abriu por dentro:
  interpretador, JIT em camadas (seção~\ref{sec:jit}), GC sobre um heap
  compartilhado (seção~\ref{sec:gc}).
\item[BEAM] A máquina virtual do Erlang (e do Elixir). O livro aberto
  \citetitle{livro:stenman} explica o nome, \termo{BEAM}{\emph{Bogdan's Erlang Abstract Machine}, hoje lido como
  \emph{Björn's Erlang Abstract Machine}: Bogumil \enquote{Bogdan} Hausman
  começou a VM na Ericsson no início dos anos~1990; Björn Gustavsson, o
  mantenedor atual, assumiu o trabalho depois.}, e a
  define como a máquina que executa o código no Erlang Runtime System.
  Das três, é a que mais se afasta da HotSpot em concorrência e falhas.
\item[GraalVM] Um nome, dois modos de execução. No modo JIT, a documentação
  diz que a GraalVM \enquote{is based on the Java HotSpot Virtual
  Machine} e entrega o código quente ao compilador Graal, escrito em Java
  (\citetitle{doc:graalvm-java}). No modo Native Image, compila tudo antes
  da execução, como no modo nativo do Quarkus
  (seção~\ref{sec:jvm-vs-nativo}).
\end{description}

Cada seção compara um eixo: concorrência, GC, falhas, carga de código e
compilação. A seção~\ref{sec:vms-tabela} junta tudo numa tabela, com a
fonte de cada célula. As versões citadas são as das páginas abertas em
18/10/2026: Erlang/OTP~29.1.2, JDK~25 e a documentação \emph{latest} da
GraalVM.

\section{Concorrência: threads, virtual threads e processos}
\label{sec:vms-concorrencia}

\paragraph{HotSpot.}
A Parte~II mostrou os dois tipos de thread do Java. A thread de plataforma
é uma thread do sistema operacional. A virtual thread
(seção~\ref{sec:virtual-threads}) é escalonada pelo próprio JDK: o
\citetitle{doc:jep444} diz que o escalonador padrão é um
\texttt{ForkJoinPool} em modo FIFO, com paralelismo igual ao número de
processadores disponíveis. Em qualquer dos dois casos, todas as threads
enxergam o mesmo heap: a especificação da JVM diz que ele é
\enquote{shared among all Java Virtual Machine threads}
(\citetitle{doc:jvms25}, §2.5.3). Por isso a Parte~II precisou falar de
\texttt{synchronized} e do modelo de memória
(seções~\ref{sec:sincronizacao} e~\ref{sec:modelo-memoria}).

\paragraph{BEAM.}
A unidade de concorrência do Erlang é o processo, que não é um processo
do sistema operacional. A documentação diz que eles são
\enquote{lightweight (grow and shrink dynamically) with a small memory
footprint} e se comunicam por sinais assíncronos; uma mensagem recebida
vai para o fim da fila do processo, a sua
\termo{mailbox}{Caixa de correio: a fila de mensagens de um processo
Erlang, lida com \texttt{receive}.} (\citetitle{doc:erlang-processes}).
Três números da documentação dão a escala:
\begin{itemize}
\item Um processo recém-criado ocupa 327 palavras de memória, das quais
  233 são o heap inicial (que inclui a pilha). O guia avisa que o número
  veio de uma execução no OTP~27 (erts-14.2.3) e pode variar entre
  versões (\citetitle{doc:erlang-eff-processes}). Em máquina de 64 bits,
  uma palavra tem 8~bytes.
\item O limite padrão de processos simultâneos de um nó é 1\,048\,576
  (opção \texttt{+P} do \citetitle{doc:erlang-erl}).
\item O número de \termo{schedulers}{Escalonadores: threads do sistema
  operacional que executam processos Erlang, cada uma com a sua fila.}
  é, por padrão, o número de processadores lógicos (opção \texttt{+S},
  mesma página).
\end{itemize}
O mesmo guia diz que tudo o que vai numa mensagem entre processos é
\emph{copiado}, com duas exceções no mesmo nó: \emph{refc binaries}
(binários grandes, contados por referência) e literais. Não há memória
compartilhada para proteger com lock: o preço é a cópia.

\paragraph{Preempção por reductions.}
Aqui está a diferença mais funda com as virtual threads. O
\citetitle{livro:stenman} descreve o escalonamento da BEAM como
\enquote{preemptive scheduling on top of cooperative scheduling}: o
processo só pode ser suspenso em certos pontos, como uma chamada de
função ou um \texttt{receive}, mas a VM conta
\termo{reductions}{Unidades de trabalho de um processo Erlang;
grosso modo, cada chamada de função gasta uma. O nome vem do Prolog, em
que cada passo de execução é uma redução.} e, quando o processo gasta a
sua cota, ele é suspenso e outro entra. O livro dá a cota: 4000
reductions (\texttt{CONTEXT\_REDS} em \texttt{erl\_vm.h}); antes do
OTP~20, eram 2000. Como em Erlang não há laço sem chamada de função, nenhum
processo segura um scheduler para sempre.

O JEP~444 diz o contrário para o Java: \enquote{The scheduler does not
currently implement time sharing for virtual threads}. Uma virtual
thread que faz conta sem bloquear fica com a carrier thread até acabar,
o mesmo alerta que o guia do Quarkus repete
(seção~\ref{sec:vt-vs-reativo}).

\paragraph{Onde a BEAM também trava.}
A preempção vale para código Erlang. Uma NIF (função nativa em C chamada
pelo Erlang) roda fora dessa contagem; o \emph{BEAM Book} recomenda
mantê-las curtas \enquote{to avoid blocking the BEAM scheduler} e usar
\emph{dirty schedulers} para operações longas. É o equivalente, na BEAM,
do \emph{pinning} das virtual threads em código nativo
(seção~\ref{sec:virtual-threads}).

\paragraph{O próprio JEP compara.}
O JEP~444 apresenta virtual threads como \emph{user-mode threads}, um
modelo que deu certo em outras linguagens, e cita \enquote{goroutines in
Go and processes in Erlang}. A comparação é justa no custo de criar e
no número de tarefas. Não é justa em dois pontos: o heap continua
compartilhado e não há preempção por tempo. Uma virtual thread é uma
thread barata, não um processo isolado.

\paragraph{GraalVM.}
No modo JIT, o modelo é o da HotSpot, em que a GraalVM se baseia
(\citetitle{doc:graalvm-java}). No Native Image, a linguagem e as APIs de
concorrência são as do Java; o que muda é a compilação e o GC, assunto
das próximas seções.

\section{Coleta de lixo: um heap ou um heap por processo}
\label{sec:vms-gc}

\paragraph{HotSpot.}
Um heap para todas as threads (§2.5.3 da \citetitle{doc:jvms25}) e um
coletor que cuida dele inteiro. A seção~\ref{sec:gc} mostrou as opções do
JDK~25: G1 por padrão em máquinas server-class (duas ou mais CPUs e
1792~MB ou mais; nas outras, Serial), ZGC para pausas curtas, Parallel e Serial para
casos específicos. Como o heap é um só, o coletor precisa lidar com objetos
de todas as threads ao mesmo tempo, e o tamanho do heap é uma decisão de
toda a aplicação.

\paragraph{BEAM.}
A documentação do coletor da BEAM descreve outra arquitetura:
\enquote{a per process generational semi-space copying collector using
Cheney's copy collection algorithm together with a global large object
space} (\citetitle{doc:erlang-gc}). Em português:
\begin{itemize}
\item cada processo tem a sua pilha e o seu heap, e o coletor trabalha
  sobre o heap de \emph{um} processo de cada vez;
\item dentro desse heap, a coleta é geracional e por cópia;
\item binários com mais de 64 bytes ficam fora dos heaps dos processos,
  num espaço global contado por referência; o processo guarda só um
  ponteiro para eles.
\end{itemize}
O \citetitle{livro:stenman} completa: quando um processo fica sem espaço
no heap, a VM primeiro tenta recuperar espaço coletando o lixo
\emph{daquele} processo. Como as mensagens são copiadas
(seção~\ref{sec:vms-concorrencia}), não há ponteiro de um heap de
processo para outro, e é isso que torna a coleta por processo possível.
O que a pesquisa encontrou ao testar heaps por thread na JVM, e o que a
HotSpot faz em vez disso (o TLAB), está na seção~\ref{sec:gc}.

\paragraph{O que isso muda na prática.}
O custo da coleta passa a ser proporcional ao processo, e um processo
com muita memória não obriga os outros a pagar por ele. A página do
coletor, porém, não diz que a coleta de um processo nunca afeta os
outros, e nenhuma das fontes abertas aqui promete isso; não leia
\enquote{GC por processo} como \enquote{nenhuma pausa global}. O lado
ruim aparece nos binários. O \emph{BEAM Book} explica que um pedaço
extraído de um binário (\emph{sub binary}) é só uma estrutura com
ponteiros para o binário original e aumenta a contagem de referências
dele: enquanto o pedaço pequeno estiver vivo, o binário grande também
estará.

\paragraph{GraalVM.}
No modo JIT, os coletores são os da HotSpot. No Native Image, a página
\citetitle{doc:graalvm-memory} lista outros: o Serial GC é o padrão,
descrito como \enquote{a simple (non-parallel, non-concurrent) stop and
copy GC}, com geração jovem e velha; o G1 existe só no Oracle GraalVM
(não na Community Edition), em Linux AMD64 e AArch64, com
\texttt{--gc=G1}; e o Epsilon não coleta nada. A
seção~\ref{sec:memoria-container} mostrou o efeito disso em container.

\paragraph{Ponte com o seu código.}
Em Java, o mais perto de \enquote{heap por processo} é \emph{não}
compartilhar estado mutável entre tarefas: objetos imutáveis
(\texttt{record}, seção~\ref{sec:records}) e mensagens copiadas entre
componentes. Isso não muda o GC da HotSpot, mas tira o estado
compartilhado do caminho, que é o que a BEAM resolve na VM.

\section{Falhas: exceções ou \emph{let it crash}}
\label{sec:vms-falhas}

\paragraph{A filosofia do Erlang.}
A tese de Joe Armstrong (\citetitle{paper:armstrong-thesis-reliable-systems},
KTH, 2003) parte de uma premissa: sistemas grandes vão para produção
\emph{com} erros, por mais que tenham sido testados. A resposta não é
evitar todo erro, e sim isolar o estrago. Na seção~2.4.3, o isolamento
tem consequências diretas: processos não compartilham nada, mensagem é o
único jeito de passar dados e, por isso, a troca de mensagens é
assíncrona. Na seção~4.3, a filosofia vira quatro lemas:
\begin{quote}
Let some other process do the error recovery. If you can't do what you
want to do, die. Let it crash. Do not program defensively.
\end{quote}

\paragraph{Exceção não é erro.}
A seção~4.4 da tese faz uma distinção que vale para Java também:
\enquote{exceptions occur when the run-time system does not know what to
do} e \enquote{errors occur when the programmer doesn't know what to do}.
Abrir um arquivo que não existe gera uma exceção, mas, se o programador
previu o caso e sabe o que fazer, não é erro: ele trata. O
\termo{let it crash}{Deixe quebrar: quando o código não sabe o que fazer,
o processo termina e outro processo cuida da recuperação.} vale para o
segundo caso, o que ninguém previu. Em vez de um \texttt{catch} genérico
que tenta seguir com estado duvidoso, o processo morre limpo.

\paragraph{Links e supervisores.}
Morrer só é seguro se alguém percebe. A documentação de processos diz que
dois processos podem ser ligados por um \emph{link}, bidirecional, e que
um processo que termina \enquote{will send exit signals to all processes
and ports that it is linked to} (\citetitle{doc:erlang-processes}). Sobre
isso o OTP monta os supervisores. Segundo \citetitle{doc:erlang-supervisor},
um supervisor é \enquote{responsible for starting, stopping, and
monitoring its child processes}, e supervisores supervisionam outros
supervisores, formando uma
\termo{supervision tree}{Árvore de supervisão: supervisores cujos filhos
são processos de trabalho ou outros supervisores.}. A página lista três
estratégias principais de reinício\footnote{A página também descreve
\texttt{simple\_one\_for\_one}, um \texttt{one\_for\_one} simplificado em
que todos os filhos são instâncias, adicionadas dinamicamente, do mesmo
processo.}:
\begin{description}
\item[\texttt{one\_for\_one}] só o filho que terminou é reiniciado;
\item[\texttt{one\_for\_all}] os outros filhos também são terminados e
  todos são reiniciados;
\item[\texttt{rest\_for\_one}] terminam e reiniciam o filho e os que
  foram iniciados depois dele.
\end{description}
E um freio: se houver mais de \texttt{intensity} reinícios em
\texttt{period} segundos, o supervisor termina todos os filhos e depois a
si mesmo, e o problema sobe um nível na árvore. Reiniciar em laço para
sempre não é recuperação.

\paragraph{Java: a exceção mata a thread, não o estado.}
A JLS (\citetitle{doc:jls25}, §11.3) diz o que acontece com uma exceção
que ninguém trata: a thread que a lançou termina, depois de executar os
\texttt{finally}, e o \emph{uncaught exception handler} é chamado. As
outras threads seguem. Como o heap é compartilhado
(seção~\ref{sec:vms-concorrencia}), os objetos que a thread estava
alterando continuam lá, visíveis para as outras, talvez pela metade. O
\enquote{morrer limpo} do Erlang depende de o processo levar o próprio
estado junto; em Java, isso só vale se você desenhar assim.

\paragraph{O que o Java tem de parecido.}
\begin{itemize}
\item \emph{Concorrência estruturada.} O \citetitle{doc:jep505} (prévia
  no JDK~25) trata um grupo de subtarefas como uma unidade, com
  tratamento de erro e cancelamento juntos, e diz que ideias como
  \enquote{Erlang's hierarchical supervisors} informaram o desenho. É
  uma API em prévia: precisa de \texttt{--enable-preview} e pode mudar.
\item \emph{Padrões de resiliência.} Timeout, retry, circuit breaker e
  bulkhead (seção~\ref{sec:resiliencia-quarkus}) limitam o estrago de uma
  dependência que falha. O bulkhead é o parente mais próximo do
  isolamento: separa recursos para que um cliente lento não leve os
  outros junto.
\item \emph{Supervisão fora da JVM.} No Kubernetes, quando o container
  falha o \emph{liveness probe} além da tolerância configurada,
  \enquote{the kubelet restarts that container}
  (\citetitle{doc:k8s-probes}). O Quarkus expõe esses probes pelo
  \citetitle{doc:quarkus-health}. É um supervisor \texttt{one\_for\_one}
  com granularidade de container, não de tarefa.
\end{itemize}
O capítulo~\ref{cap:decisoes-vm} volta a essas pontes e às diferenças
que ficam.

\section{Trocar código com o sistema rodando}
\label{sec:vms-hot-code}

\paragraph{BEAM: duas versões de cada módulo.}
O Erlang foi feito para trocar código sem parar o sistema, o
\termo{hot code loading}{Carga de código a quente: substituir um módulo
enquanto o sistema roda, sem reiniciar.}. A página
\citetitle{doc:erlang-code-loading} descreve a regra:
\begin{itemize}
\item um módulo pode existir em duas variantes ao mesmo tempo, a
  \emph{current} e a \emph{old}; carregar uma versão nova torna a anterior
  \emph{old};
\item \enquote{Fully qualified function calls always refer to current
  code}: uma chamada \texttt{modulo:funcao()} sempre vai para a versão
  nova, e é assim que um processo migra;
\item ao carregar uma terceira versão, o code server remove a mais antiga
  \enquote{and any processes lingering in it are terminated}.
\end{itemize}
O \citetitle{doc:beam-book} tira a consequência para quem escreve o
código: um laço de servidor que não faz nenhuma chamada qualificada
continua rodando o código velho. A troca a quente funciona porque o
código foi escrito para ela.

\paragraph{HotSpot: redefinição com limites.}
A JVM também troca código em execução, mas bem menos. A
\citetitle{doc:jvmti25} diz que a redefinição de uma classe pode mudar
corpos de métodos, o constant pool e atributos, e \enquote{must not add,
remove or rename fields or methods, change the signatures of methods,
change modifiers, or change inheritance}. A documentação de
\citetitle{doc:instrumentation25} completa: um método redefinido que tem
frames ativos continua executando o bytecode antigo nesses frames; a
versão nova vale para as próximas chamadas. É a mesma ideia das duas
versões do Erlang, restrita ao corpo dos métodos. É o que depuradores e
agentes de instrumentação usam, e não um mecanismo de deploy.

\paragraph{Quarkus: reinício rápido no lugar de troca a quente.}
O dev mode do Quarkus resolve o problema de outro jeito. O
\citetitle{doc:quarkus-class-loading} separa um class loader persistente
(o \emph{base runtime}, que carrega o que não muda) de um \emph{runtime
class loader} com as classes da aplicação, que \enquote{is recreated when
the application is restarted}. Em vez de redefinir classes, o Quarkus
descarta o class loader e cria outro: na prática, reinicia só a parte que
mudou. A mesma página avisa que isso não tem efeito em builds de produção.

\paragraph{Native Image: nada novo depois do build.}
No outro extremo está o executável nativo. A
\termo{closed-world assumption}{Suposição de mundo fechado: todo o código
que pode executar é conhecido no build (seção~\ref{sec:jvm-vs-nativo}).}
diz que, construída a imagem, \enquote{no new elements can be added at
run time, for example, through class loading}
(\citetitle{doc:graalvm-native-basics}). Trocar código significa gerar
outro binário e trocar o processo.

\paragraph{A decisão que isso muda.}
Em Java, o \enquote{hot code loading} de produção é o deploy sem
downtime: subir réplicas novas, mover o tráfego e desligar as velhas. O
problema que o Erlang resolve dentro da VM (duas versões convivendo) você
resolve fora dela, e com o mesmo cuidado: durante o rollout, versões
diferentes do serviço atendem ao mesmo tempo, então contratos de API e
esquema de banco precisam aceitar as duas (é para isso que servem os
testes de contrato, como em \citetitle{artigo:consumer-driven-contracts}).

\section{Compilação: JIT especulativo, JIT de carga e AOT}
\label{sec:vms-jit-aot}

\paragraph{HotSpot: especular com perfil.}
A seção~\ref{sec:jit} mostrou o ciclo: interpretar, medir, compilar o que
é quente com C1 e depois C2, especular sobre o que o perfil mostrou e,
se a especulação falha, \emph{desotimizar} e voltar ao interpretador.
O capítulo~\ref{cap:historia-vms} mostra de onde veio esse ciclo.

\paragraph{Graal como JIT.}
A página \citetitle{doc:graalvm-compiler} descreve o Graal como
\enquote{a dynamic compiler, written in Java, that transforms bytecode
into machine code}, integrado à HotSpot. Entre as fases de otimização,
cita \emph{inlining} agressivo, \emph{inlining} polimórfico e uma
\emph{partial escape analysis} que \enquote{can remove the costly
allocations of certain objects}: a escape analysis da seção~\ref{sec:jit},
aplicada também quando o objeto só escapa em alguns caminhos. O
ciclo de especulação e desotimização é o mesmo da HotSpot; muda o
compilador do último nível.

\paragraph{Truffle: o interpretador vira compilador.}
O Graal também é a base do Truffle, que a documentação apresenta como
uma biblioteca para implementar linguagens que \enquote{automatically
deriv[es] high-performance code from interpreters}
(\citetitle{doc:truffle}). O mecanismo é a
\termo{partial evaluation}{Avaliação parcial: especializar um programa
para parte das suas entradas já conhecidas. Aqui, especializar o
interpretador para um programa fixo, o que produz código compilado desse
programa.}: \textcite{paper:truffle-partial-evaluation} definem a
linguagem só escrevendo um interpretador; o código compilado sai
automaticamente por avaliação parcial e, quando uma especulação falha,
\enquote{execution transfers back to the interpreter}. É a desotimização
da HotSpot generalizada para qualquer linguagem.

\paragraph{Native Image: compilar tudo antes.}
\textcite{paper:native-image-initialize-once}, o paper do Native Image,
descrevem o processo: uma aplicação iterativa de análise de
\emph{points-to} e de \emph{snapshot} do heap, seguida de compilação AOT
com um compilador otimizador. Código de inicialização pode rodar no
build, e a execução começa com um heap já populado, o \emph{image heap}
(seção~\ref{sec:jvm-vs-nativo}); é por isso que o trabalho que o Quarkus
adianta para o build (capítulo~\ref{cap:build-time}) rende tanto no nativo. Sem perfil de
execução e sem class loading, não há o que especular em runtime; a
seção~\ref{sec:jvm-vs-nativo} mostrou o efeito no throughput de pico.

\paragraph{BEAM: JIT sem especulação.}
O JIT da BEAM, o BeamAsm, faz outra coisa. A documentação diz que ele
\enquote{provides load-time conversion of Erlang BEAM instructions into
native code on x86-64 and aarch64}: traduz cada módulo quando ele é
carregado, eliminando o custo de despacho do interpretador e
especializando cada instrução pelos tipos dos argumentos. E avisa:
\enquote{BeamAsm does hardly any cross instruction optimizations}
(\citetitle{doc:erlang-beamasm}). Não há perfil, não há inlining
especulativo, não há desotimização. Em troca, o JIT é simples, o código
nativo existe desde a carga e a troca a quente de módulos
(seção~\ref{sec:vms-hot-code}) continua funcionando como antes.

\paragraph{Três respostas para \enquote{quando compilar}.}
\begin{description}
\item[Durante a execução, com perfil] HotSpot e GraalVM JIT: demoram a
  aquecer, chegam ao melhor código de pico, dependem de especulação e
  desotimização.
\item[Na carga, sem perfil] BeamAsm: o código nativo existe desde a
  carga do módulo, com otimização local por instrução.
\item[No build] Native Image: startup rápido e memória menor, ao custo
  do mundo fechado e de compilar sem ver a execução real.
\end{description}

\section{Tabela-resumo}
\label{sec:vms-tabela}

A tabela junta as seções anteriores. Cada célula traz a fonte; onde não
encontramos número publicado e comparável, a célula diz isso em vez de
inventar um. Números de desempenho são do fornecedor e servem como ordem
de grandeza (capítulo~\ref{cap:benchmarks}).

\footnotesize
\begin{longtable}{@{}>{\raggedright\arraybackslash}p{1.9cm} >{\raggedright\arraybackslash}p{3.4cm} >{\raggedright\arraybackslash}p{3.4cm} >{\raggedright\arraybackslash}p{3.6cm}@{}}
\toprule
Eixo & HotSpot (JDK~25) & BEAM (OTP~29.1.2) & GraalVM (JIT e Native
  Image) \\
\midrule \endhead
Modelo de concorrência & Threads de plataforma (do SO) e virtual
  threads num \texttt{ForkJoinPool}, sem \emph{time sharing}; heap
  compartilhado (\citetitle{doc:jep444}; \citetitle{doc:jvms25},
  §2.5.3) & Processos leves, sem memória compartilhada, mensagens
  copiadas para a mailbox; preempção por reductions
  (\citetitle{doc:erlang-processes};
  \citetitle{doc:erlang-eff-processes}; \citetitle{livro:stenman}) &
  O da HotSpot, na qual a GraalVM se baseia
  (\citetitle{doc:graalvm-java}); no Native Image, as mesmas APIs Java
  \\
\addlinespace
GC & Um heap para todas as threads; G1 padrão em máquinas
  server-class (Serial nas outras), ZGC, Parallel
  (seção~\ref{sec:gc}; \citetitle{doc:gc-collectors}) & Por processo,
  geracional por cópia (Cheney), mais espaço global de binários
  contados por referência (\citetitle{doc:erlang-gc}) & JIT: os da
  HotSpot. Native: Serial (padrão, \emph{stop and copy}); G1 só no
  Oracle GraalVM em Linux; Epsilon (\citetitle{doc:graalvm-memory}) \\
\addlinespace
Tratamento de falha & Exceção não tratada termina a thread; o estado no
  heap fica (\citetitle{doc:jls25}, §11.3); resiliência por bibliotecas
  (seção~\ref{sec:resiliencia-quarkus}) & \emph{Let it crash}: links,
  sinais de saída e árvores de supervisão com limite de reinícios
  (\citetitle{paper:armstrong-thesis-reliable-systems};
  \citetitle{doc:erlang-supervisor}) & Igual à HotSpot nos dois modos:
  é a mesma linguagem Java (\citetitle{doc:jls25}, §11.3;
  \citetitle{doc:graalvm-java}) \\
\addlinespace
Carga de código & Class loading dinâmico
  (seção~\ref{sec:class-loading}); redefinição só de corpos de métodos
  (\citetitle{doc:jvmti25}) & Duas versões por módulo; chamada
  qualificada migra para a nova; a terceira carga mata quem ficou na
  mais antiga (\citetitle{doc:erlang-code-loading}) & JIT: como a
  HotSpot. Native: mundo fechado, nada novo em runtime
  (\citetitle{doc:graalvm-native-basics}) \\
\addlinespace
Compilação & Interpretador + C1/C2 com perfil, especulação e
  desotimização (seção~\ref{sec:jit}; \citetitle{doc:hotspot-perf}) &
  BeamAsm: tradução na carga, sem otimização entre instruções
  (\citetitle{doc:erlang-beamasm}) & JIT: Graal, escrito em Java, com
  \emph{partial escape analysis} (\citetitle{doc:graalvm-compiler}).
  Native: AOT após análise de \emph{points-to}
  (\citetitle{paper:native-image-initialize-once}) \\
\addlinespace
Startup e memória & Quarkus em modo JVM: 4417~ms até a 1ª requisição,
  304~MiB de RSS. Perf-lab do Quarkus, 21/04/2026, Quarkus~3.34.3,
  JDK~25.0.2 (fornecedor; seção~\ref{sec:jvm-vs-nativo}) & 327
  palavras por processo novo (OTP~27, \citetitle{doc:erlang-eff-processes}).
  Não achamos medição publicada de startup e RSS de um serviço
  comparável ao do Quarkus & Native: 581~ms e 95~MiB na mesma medição
  do Quarkus (fornecedor; seção~\ref{sec:jvm-vs-nativo}). O paper do
  Native Image fala em startup \enquote{up to two orders of magnitude}
  melhor que a HotSpot (\citetitle{paper:native-image-initialize-once};
  Oracle Labs, fornecedor, 2019) \\
\bottomrule
\end{longtable}
\normalsize

\paragraph{Como ler a tabela.}
A linha de startup é a única com números, e eles não se comparam entre
colunas: a medição do Quarkus é de um serviço REST com banco, e as 327
palavras da BEAM são o custo de \emph{um} processo, não de um serviço.
O que a tabela mostra bem são as escolhas de projeto: a HotSpot aposta em
heap compartilhado e JIT especulativo; a BEAM, em isolamento e troca a
quente; o Native Image, em saber tudo no build.

\section{Onde aprofundar}

\subsection*{Livros e teses}
\small
\begin{longtable}{@{}p{5.2cm} p{7cm}@{}}
\toprule
Fonte & Para quê \\
\midrule \endhead
\citetitle{livro:stenman} (Stenman, aberto, CC BY 4.0) & A BEAM por
  dentro: processos, escalonamento por reductions, GC, carga de código \\
\citetitle{paper:armstrong-thesis-reliable-systems} (Armstrong, KTH,
  2003) & A filosofia do Erlang/OTP; leia os capítulos~2 e~4 \\
\citetitle{livro:cesarini2016} (Cesarini e Vinoski, 2016) & Supervisão,
  \emph{behaviours} e confiabilidade em sistemas OTP reais \\
\citetitle{paper:history-of-erlang} (Armstrong, HOPL~III, 2007) & A
  história do Erlang contada pelo criador \\
\bottomrule
\end{longtable}
\normalsize

\subsection*{Documentação}
\small
\begin{longtable}{@{}p{5.2cm} p{7cm}@{}}
\toprule
Fonte & Para quê \\
\midrule \endhead
\citetitle{doc:erlang-processes}, \citetitle{doc:erlang-eff-processes}
  & Processos, sinais, links e custo de memória \\
\citetitle{doc:erlang-gc} & O coletor por processo e os binários
  grandes \\
\citetitle{doc:erlang-supervisor} & Estratégias de reinício e limite de
  intensidade \\
\citetitle{doc:erlang-code-loading} & As duas versões de um módulo \\
\citetitle{doc:erlang-beamasm} & Como funciona o JIT da BEAM \\
\citetitle{doc:graalvm-compiler}, \citetitle{doc:truffle} & O Graal como
  JIT e como base do Truffle \\
\citetitle{doc:jep505} & Concorrência estruturada no Java, inspirada em
  parte nos supervisores do Erlang \\
\bottomrule
\end{longtable}
\normalsize

\subsection*{Papers}
\textcite{paper:truffle-partial-evaluation} e
\textcite{paper:native-image-initialize-once} explicam, com dados, os
dois modos da GraalVM que este capítulo resumiu.

\section{Desafios práticos}

\desafio{Crie um milhão de virtual threads que dormem um segundo cada e
meça o heap com JFR. Depois, num shell Erlang, crie um milhão de
processos que esperam uma mensagem e veja a memória com
\texttt{erlang:memory()}. O que cada número inclui? Por que comparar os
dois diretamente é enganoso?}

\desafio{Escreva uma tarefa que faz conta sem bloquear (um laço de CPU)
e rode, ao mesmo tempo, mais tarefas desse tipo do que núcleos, em
virtual threads. Meça a latência de uma tarefa curta que chega no meio.
Explique o resultado com a frase do JEP~444 sobre \emph{time sharing} e
compare com o que a BEAM faria com reductions.}

\desafio{Escreva um serviço Java em que uma thread altera dois campos de
um objeto compartilhado e lança uma exceção entre as duas alterações.
O que as outras threads veem? Redesenhe para que a falha não deixe estado
pela metade, sem usar \texttt{synchronized}.}

\desafio{Com um agente Java ou o depurador da IDE, redefina o corpo de um
método durante a execução. Depois tente acrescentar um campo à classe.
Qual erro aparece, e qual regra da JVM TI ele aplica?}

\desafio{Monte, em Erlang ou Elixir, um supervisor \texttt{one\_for\_one}
com um filho que quebra a cada mensagem. Encontre os valores de
\texttt{intensity} e \texttt{period} que fazem o supervisor desistir.
Qual seria o equivalente no seu serviço Quarkus: o circuit breaker, o
\emph{liveness probe} ou nenhum dos dois?}

